\section{异步执行与错误检查}

\subsection{核函数启动默认是异步的}

\par\noindent\begin{minipage}{\linewidth}
一次核函数启动通常先把工作提交给 GPU，然后让主机线程继续执行。此行为称为\term{异步调用}{asynchronous call}。因此，下面的 CPU 计时方式主要测到“提交请求”的时间，不一定覆盖 GPU 完成计算所需的时间：

\begin{lstlisting}[style=cuda,numbers=none]
start_timer();
vecAddKernel<<<B, T>>>(A_d, B_d, C_d, n);
stop_timer();      // usually too early
\end{lstlisting}
\end{minipage}\par

\par\noindent\begin{minipage}{\linewidth}
要让主机等待此前提交的设备工作完成，可调用
\begin{lstlisting}[style=cuda,numbers=none]
cudaDeviceSynchronize();
\end{lstlisting}
\end{minipage}\par
该函数提供主机与设备之间的\term{同步点}{synchronization point}，并返回此前异步工作中暴露的错误状态。

\slidefig{0.94}{page-32.png}{异步核函数调用与 \code{cudaDeviceSynchronize()}}

当核函数后紧跟设备到主机的 \code{cudaMemcpy} 时，复制不能在产生结果的核函数完成之前得到最终输出。更精确的性能测量通常使用 CUDA\term{事件}{event}，并在需要时配合\term{流同步}{stream synchronization}或\term{性能分析器}{profiler}。

\subsection{两类错误：启动错误与执行错误}

绝大多数 CUDA 运行时 API 都返回 \code{cudaError\_t}。\code{cudaSuccess} 表示成功。其他错误码可交给 \code{cudaGetErrorString}，转换为可读消息。

核函数调用本身没有普通返回值，因此常分两步检查：
\begin{enumerate}
  \item \code{cudaGetLastError()}：检查非法启动配置等立即可见的\term{启动错误}{launch error}；
  \item \code{cudaDeviceSynchronize()}：等待执行完成，并暴露越界访问等\term{异步执行错误}{asynchronous execution error}。
\end{enumerate}

\par\noindent\begin{minipage}{\linewidth}
\begin{lstlisting}[style=cuda,numbers=none]
vecAddKernel<<<numBlocks, threadsPerBlock>>>(A_d, B_d, C_d, n);
CUDA_CHECK(cudaGetLastError());
CUDA_CHECK(cudaDeviceSynchronize());
\end{lstlisting}
\end{minipage}\par

\begin{warningbox}[title={简化示例与工程代码的区别}]
用于突出概念的短代码通常会省略错误检查，也可能包含简化写法。工程代码应检查内存分配、数据复制、核函数启动与核函数执行的返回状态。
\end{warningbox}
